\documentclass[runningheads]{llncs}
\usepackage{fontspec}
\usepackage{xeCJK}
\setmainfont{TeX Gyre Termes}
\setCJKmainfont[AutoFakeSlant=true]{FandolSong-Regular}
\setCJKsansfont{FandolHei-Regular}
\setCJKmonofont{FandolFang-Regular}
\usepackage{amsmath,amssymb,booktabs}
\usepackage{algorithm}
\usepackage{algpseudocode}
\usepackage[hidelinks]{hyperref}
\usepackage{tikz}
\usetikzlibrary{arrows.meta,positioning,calc}

\title{一种王者荣耀拆塔速推流的可行性分析与实现}
\titlerunning{王者荣耀事件驱动拆塔速推流}

\author{杜宇龙}
\authorrunning{杜宇龙}

\institute{娱乐策略研究组\\
\email{contact@example.com}}

\begin{document}
\maketitle

\begin{abstract}
本文从事件驱动系统视角研究王者荣耀中的拆塔速推流。与常见“按分路重复循环推进”的思路不同，本文提出“节奏积分账本”与“分路分配有限状态机”组合框架，并由轻量调度器在兵线同步、视野突增、目标计时与撤退安全变化时触发动作。该方法强调在崩盘概率受控前提下完成结构性收益转化。受控自定义对局观察表明，该框架可提高前两座外塔转化稳定性，并降低爆发推进失败后的波动。
\keywords{王者荣耀 \and 拆塔速推 \and 事件驱动宏观控制 \and 有限状态机 \and 团队口令协议}
\end{abstract}

\section{引言}
社区语境中的“速推”常被简化为短口号，实际执行却容易在信息噪声下失真。多数失败并非源于单次操作，而是三类条件错位：兵线几何、守方显露模式、撤退通道质量。
因此，本文将拆塔速推建模为事件驱动控制问题：维护结构推进积分，并且仅在“收益门槛”和“安全门槛”同时满足时发起推进动作。

\section{运行假设与遥测信号}
在每个决策时刻 $t$，系统采样以下信号：
\begin{itemize}
\item 兵线同步得分 $W_t$（中边路波次协同程度）；
\item 守方可见性向量 $V_t$（敌方近期显位信息）；
\item 中立目标计时桶 $O_t$（暴君/主宰窗口）；
\item 逃生就绪度 $E_t$（位移技能与撤退路径可用性）。
\end{itemize}
状态统一表示为 $S_t=(W_t,V_t,O_t,E_t)$，并在每个沟通周期增量更新。

\section{节奏积分账本模型}
本文不直接求解长时域最优路径，而是维护可再生标量资源“节奏积分” $B_t$：
\[
B_{t+1}=\min\!\left(B_{\max},\;
\eta B_t
+\xi\cdot \mathrm{WaveSync}_t
+\omega\cdot \mathrm{VisionEdge}_t
-\chi\cdot \mathrm{SkirmishDrag}_t\right),
\]
其中 $\eta\in(0,1]$ 为记忆衰减因子。

解释如下：
\begin{itemize}
\item 兵线同步与视野优势会为后续推进充能；
\item 低收益拉扯会消耗积分；
\item 高强度拆塔动作必须在积分达标后才允许触发。
\end{itemize}

同时定义崩盘风险估计：
\[
H_t=\sigma(\rho_0+\rho_1F_t+\rho_2C_t-\rho_3R_t),
\]
其中 $F_t$ 为战争迷雾不确定性，$C_t$ 为关键技能冷却缺口，$R_t$ 为撤退路径质量。

\begin{table}[t]
\centering
\caption{事件驱动控制器的核心标量}
\begin{tabular}{ll}
\toprule
信号 & 含义 \\
\midrule
$B_t$ & 当前可用于结构转化的节奏积分 \\
$H_t$ & 崩盘风险估计（0 到 1） \\
$W_t$ & 分路兵线同步得分 \\
$V_t$ & 守方可见性质量 \\
$E_t$ & 逃生准备度 \\
\bottomrule
\end{tabular}
\end{table}

\section{分路分配有限状态机}
宏观推进由四个状态组成：
\begin{enumerate}
\item \textbf{STACK}：堆线与阵型重置；
\item \textbf{DRAG}：利用目标姿态或假动作牵扯守方；
\item \textbf{MELT}：短时拆塔爆发；
\item \textbf{SCREEN}：封堵侧翼并确保撤退走廊。
\end{enumerate}

\begin{figure}[t]
\centering
\begin{tikzpicture}[
  >=Latex,
  node distance=10mm and 16mm,
  state/.style={draw, rounded corners, minimum width=22mm, minimum height=8mm, align=center, font=\small}
]
\node[state] (stack) {STACK};
\node[state, right=of stack] (drag) {DRAG};
\node[state, right=of drag] (melt) {MELT};
\node[state, below=of drag] (screen) {SCREEN};

\draw[->] (stack) -- node[above, font=\scriptsize] {$W_t\uparrow,\ V_t\uparrow$} (drag);
\draw[->] (drag) -- node[above, font=\scriptsize] {$B_t\ge B_{\text{melt}}$} (melt);
\draw[->] (melt) -- node[right, font=\scriptsize] {守方回补} (screen);
\draw[->] (screen) -- node[below, font=\scriptsize] {重置完成} (stack);
\draw[->,bend left=22] (drag.north) to node[above, font=\scriptsize, fill=white, inner sep=1pt] {$H_t>H_{\max}$} (stack.north);
\draw[->,bend right=24] (melt.south) to node[right, font=\scriptsize, fill=white, inner sep=1pt] {积分不足} (screen.east);
\end{tikzpicture}
\caption{分路分配有限状态机。}
\label{fig:hok-fsa-zh}
\end{figure}

\section{口令协议}
团队沟通层采用五个互斥口令，避免语义重叠：
\begin{table}[t]
\centering
\caption{事件驱动拆塔流口令语义}
\begin{tabular}{p{1.7cm}p{3.1cm}p{5.6cm}}
\toprule
口令 & 触发条件 & 强制动作 \\
\midrule
\texttt{STACK} & 积分不足或兵线不同步 & 控接触深度，堆下一波线，保持关键技能对等。 \\
\texttt{DRAG} & 中立目标窗口开启且可见性占优 & 通过目标姿态或假推进拉扯守方站位。 \\
\texttt{MELT} & $B_t\!\ge\!B_{\text{melt}}$ 且 $H_t\!\le\!H_{\max}$ & 发起短时拆塔爆发，严格执行停线。 \\
\texttt{SCREEN} & 爆发后检测到回补或侧翼风险 & 封堵回防通道并保护撤退路线。 \\
\texttt{RESET} & 风险突增或口令不同步 & 立即撤出，清理节奏债务并返回 STACK。 \\
\bottomrule
\end{tabular}
\end{table}

\section{事件驱动调度器}
调度器处理事件集合
\[
\mathcal{E}=\{\mathrm{WAVE\_SYNC},\mathrm{VIS\_SPIKE},\mathrm{OBJ\_OPEN},
\mathrm{DANGER\_SPIKE},\mathrm{COOLDOWN\_READY}\}.
\]
与时域最大化策略不同，本文方案在每个事件边界即时发令。

\begin{algorithm}[H]
\caption{拆塔速推事件调度流程}
\begin{algorithmic}[1]
\Require 状态 $S_t$，节奏积分 $B_t$，风险估计 $H_t$
\Ensure 口令 \{\texttt{STACK}, \texttt{DRAG}, \texttt{MELT}, \texttt{SCREEN}, \texttt{RESET}\}
\If{$H_t > H_{\max}$}
  \State \Return \texttt{RESET}
\EndIf
\If{事件为 \texttt{WAVE\_SYNC} 且 $B_t < B_{\text{melt}}$}
  \State \Return \texttt{STACK}
\EndIf
\If{事件为 \texttt{OBJ\_OPEN} 且可见性优势为正}
  \State \Return \texttt{DRAG}
\EndIf
\If{$B_t \ge B_{\text{melt}}$ 且 $H_t \le H_{\max}$ 且逃生就绪}
  \State \Return \texttt{MELT}
\EndIf
\If{爆发结束或侧翼风险上升}
  \State \Return \texttt{SCREEN}
\EndIf
\State \Return \texttt{STACK}
\end{algorithmic}
\end{algorithm}

\begin{figure}[t]
\centering
\begin{tikzpicture}[x=1.0cm,y=0.75cm,>=Latex]
\draw[->] (0,0) -- (10.5,0) node[right] {时间};
\draw (0,1.4) node[left, align=right] {\scriptsize 积分\\\scriptsize 窗口};
\draw[thick,blue!70] (0.2,1.1) -- (2.4,1.1) -- (2.4,1.7) -- (5.3,1.7) -- (5.3,1.0) -- (8.8,1.0);
\draw[dashed] (5.3,0) -- (5.3,2.1) node[above, font=\scriptsize] {MELT 爆发};
\draw[dashed] (7.4,0) -- (7.4,2.1) node[above, font=\scriptsize] {SCREEN};
\draw[thick,teal!80!black] (0.8,0.45) -- (4.9,0.45) node[midway, above, font=\scriptsize] {STACK/DRAG};
\draw[thick,orange!85!black] (5.4,0.45) -- (7.3,0.45) node[midway, above, font=\scriptsize] {MELT};
\draw[thick,purple!80!black] (7.5,0.45) -- (9.4,0.45) node[midway, above, font=\scriptsize] {SCREEN/RESET};
\end{tikzpicture}
\caption{事件相位与积分消耗时序示意。}
\label{fig:hok-timeline-zh}
\end{figure}

\section{受控评估}
我们将本文调度器与“团战优先”宏观基线在受控自定义对局中比较。观测到的趋势包括：
\begin{itemize}
\item 前两座外塔转化稳定性更高；
\item 拆塔爆发后的崩盘事件更少；
\item 推进失败后因 \texttt{RESET} 分支而恢复更快。
\end{itemize}
收益主要体现在“信息不完全但可争夺”的中盘回合。

\section{失效场景与实践注意}
若队伍违反口令纪律，尤其在 \texttt{MELT} 后继续无条件追击，系统优势会快速衰减。另一个失效点是可见性信号陈旧：当 $V_t$ 滞后真实战场时，\texttt{DRAG} 容易高估牵扯效果并诱发提前崩盘。

\section{结论}
本文表明，王者荣耀拆塔速推流可被实现为“节奏积分账本 + 有限状态机 + 事件调度器”的工程化系统。该表述不同于固定循环推进脚本，并在娱乐向实战场景中提供更清晰的失效处理能力。

\paragraph{致谢.}
本文档用于个人主页模板演示与娱乐化策略写作场景。

\begin{thebibliography}{3}
\bibitem{hok-note}
腾讯游戏：王者荣耀赛季平衡性更新说明（2026）

\bibitem{macro-log}
娱乐策略研究组：自定义宏观对局日志（2026）

\bibitem{event-control}
杜宇龙：MOBA 结构推进的事件驱动宏观调度方法，工作笔记（2026）
\end{thebibliography}

\end{document}
